%----------------------------------------------------------------------
%%% RESULTS
%----------------------------------------------------------------------
% !TEX root = ./Main.tex


We are now in a position to state and discuss our main results.
% Following the equilibrium classification scheme that we presented in \cref{sec:preliminaries},
For convenience, we will present our results in order of increasing structure, starting with stable policies, and then moving on to \acl{SOS} and deterministic Nash policies.
All proofs are deferred to the appendix.


%----------------------------------------------------------------------
%% Stable
%----------------------------------------------------------------------
\subsection{Asymptotic convergence to stable Nash policies}

Our first convergence result concerns Nash policies that satisfy the stability requirement $\braket{\gvalue(\policy)}{\policy-\eq} < 0$ of \cref{def:refine}.
In this case, we have the following guarantee:

\begin{restatable}{theorem}{stable}
\label{thm:stable}
Let $\eq$ be a stable Nash policy, and let $\curr$ be the sequence of play generated by \eqref{eq:PG} with step-size $\curr[\step] = \step/(\run + \runoff)^{\pexp}$, $\pexp \in (1/2,1]$,
and policy gradient estimates such that $\pexp + \bexp >1$ and $\pexp - \sexp > 1/2$ as per \eqref{eq:schedule}.
Then there exists a neighborhood $\nhd$ of $\eq$ in $\policies$ such that, for any given $\conf>0$, we have
\begin{equation}
\label{eq:primal-local}
\probof{\text{$\curr$ converges to $\eq$} \given \init\in\nhd}
	\geq 1-\conf
\end{equation}
provided that $\step$ is small enough \textpar{or $\runoff$ large enough} relative to $\conf$.
\end{restatable}

\begin{restatable}{corollary}{corstable}
\label{cor:stable}
Suppose that \crefrange{mod:full}{mod:value} are run with a step-size of the form $\curr[\step] = \step/(\run+\runoff)^{\pexp}$, $\pexp>1/2$, and if applicable,
%an exploration parameter $\curr[\mix] = \mix/\run^{\pexp/2}$.
an exploration parameter $\curr[\mix] = \mix/(\run+\runoff)^{\rexp}$ such that $1-\pexp < \rexp < 2\pexp-1$.
Then:
\begin{itemize}
\item
For \cref{mod:full,mod:stoch}:
the conclusions of \cref{thm:stable} hold as stated.
\item
For \cref{mod:value}:
the conclusions of \cref{thm:stable} hold as long as $\pexp>2/3$.
\end{itemize}
\end{restatable}


We note here that \cref{thm:stable}
% \textendash\ and thus, by extension, \cref{cor:stable} \textendash\
provides a trajectory convergence guarantee which is otherwise quite difficult to obtain even in structured stochastic games.
For example, if we zoom in on the class of stochastic potential (or min-max) games, the existing guarantees in the literature concern the ``best iterate'' of the algorithm, \cf \cite{LOPP22,zhang2021gradient} and references therein.
Because of this, said guarantees do not apply to the actual trajectory of play generated by \eqref{eq:PG};
this makes them less suitable for agent-based learning where the players involved are learning ``as they go'', as opposed to \emph{simulating} the game in order to approximately compute an equilibrium policy offline.

We should also note that the convergence guarantees of \cref{thm:stable} hold locally with arbitrarily high probability.
Without further assumptions, it is not possible to obtain global trajectory convergence guarantees that hold with probability $1$, even in single-state games \textendash\ that is, the case of learning in finite normal form games.
%In this (much simpler) setting, the well-known impossibility result of \citet{HMC00,HMC03} shows that it is not possible to expect convergence to \acl{NE} in all games \textendash\ not even locally.
%In this regard, the local convergence caveat in \cref{thm:stable} cannot be lifted without further structural properties in place \textendash\ such as the existence of a potential function in the spirit of \cite{LOPP22}.
\beginrev
The reason for this locality is twofold:
First, equilibrium policies are not unique in general, and gradient-based dynamics may also admit non-equilibrium attractors, such as limit cycles and the like \cite{HMC21,MPP18,MLZF+19,MHC22}.
As a result, in the presence of multiple equilibria/attractors, the best one can hope for is a local equilibrium convergence result, conditioned on the basin of attraction of said equilibrium (as per \cref{thm:stable}).

The second obstruction to a global, unconditional convergence result is probabilistic in nature, and has to do with the randomness that enters the learning process (\eg in the estimation of policy gradients via the \reinforce).
In this case, no matter how close one starts to an equilibrium policy, there is always a finite, non-zero probability that an unlucky realization of the noise can drive the process away from its basin, possibly never to return.
This issue can only be overcome in games where $\policies$ is partitioned (up to a set of measure zero) into basins of attraction of equilibrium policies.
However, this can only occur in games with a sufficiently strong global structure, like potential stochastic games, two-player zero-sum games and the like;
in complete generality, locality cannot be lifted, even in single-state problems \cite{FVGL+20,GVM21}.
\endedit


%----------------------------------------------------------------------
%% SOS
%----------------------------------------------------------------------
\subsection{Convergence to \acl{SOS} policies}

Albeit valuable as an asymptotic convergence guarantee, \cref{thm:stable} does not provide an indication of how long it will take players to actually converge to a Nash policy.
Of course, in full generality, it is not plausible to expect to be able to derive such a convergence rate because the stability requirement provides no indication on how fast the players' policy gradients stabilize near a solution.
This kind of estimate is provided by the second-order sufficient condition \eqref{eq:SOS}, which allows us to establish sufficient control over the sequence of play as indicated by the following theorem.


\begin{restatable}{theorem}{SOS}
\label{thm:SOS}
Let $\eq$ be a \acl{SOS} policy,
let $\basin$ be a neighborhood of $\eq$ such that \eqref{eq:strong} holds on $\basin$,
and
let $\curr$ be the sequence of play generated by \eqref{eq:PG} with step-size $\curr[\step] = \step/(\run + \runoff)^{\pexp}$, $\pexp \in (1/2,1]$, and policy gradient estimates such that $\pexp + \bexp > 1$ and $\pexp - \sexp > 1/2$ as per \eqref{eq:schedule}.
Then:
\begin{enumerate}
\item
There exists a neighborhood $\nhd$ of $\eq$ in $\policies$ such that, for any confidence level $\conf>0$, the event
\begin{equation}
\label{eq:good}
\event
	= \braces{\text{$\curr \in \basin$ for all $\run=\running$}}
\end{equation}
occurs with probability
%\begin{equation}
%\label{eq:good-prob}
\(
\probof{\event \given \init\in\nhd} \geq 1-\conf
\)
%\end{equation}
%if $\init\in\init[\nhd]$ and
if $\runoff$ is large enough relative to $\conf$.
\item
The sequence $\curr$ converges to $\eq$ \acl{wp1} on $\event$;
in particular, we have
\begin{equation}
\label{eq:prob-converge}
\probof{\text{$\curr$ converges to $\eq$} \given \init\in\nhd}
	\geq 1-\conf
\end{equation}
if $\runoff$ is large relative to $\conf$.
Moreover, conditioned on $\event$ and taking $\qexp = \min\{\bexp,\pexp - 2\sexp\}$, we have
\begin{equation}
\label{eq:conv-rate}
\exof{\norm{\curr - \eq}^{2} \given \event}
	=
	\begin{cases}
	\bigoh(1/\run^{2\strong\step})
		&\quad
		\text{if $\pexp=1$ \, and \, $2\strong\step<\qexp$},
		\\
	\bigoh(1/\run^{\qexp})
		&\quad
		\text{otherwise}.
%		\\
%	\bigoh(1/\run^{\qexp})
%		&\quad
%		\text{if $\pexp<1$},
	\end{cases}
\end{equation}
%where $\qexp = \min\{\bexp,\pexp - 2\sexp\}$.
\end{enumerate}
\end{restatable}

\begin{restatable}{corollary}{corSOS}
\label{cor:SOS}
Suppose that \crefrange{mod:full}{mod:value} are run with a step-size of the form $\curr[\step] = \step/(\run+\runoff)^{\pexp}$, $\pexp>1/2$, and if applicable, an exploration parameter $\curr[\mix] = \mix/(\run+\runoff)^{\pexp/2}$.
Then:
\begin{itemize}
\item
For \cref{mod:full,mod:stoch}:
the conclusions of \cref{thm:SOS} hold with $\qexp=\pexp$;
in particular, \eqref{eq:conv-rate} gives an $\bigoh(1/\run)$ rate of convergence if $\pexp=1$ and $2\strong\step > \qexp$.
\item
For \cref{mod:value}:
the conclusions of \cref{thm:SOS} hold for $\pexp>2/3$ with $\qexp = \pexp/2$;
in particular, \eqref{eq:conv-rate} gives an $\bigoh(1/\sqrt{\run})$ rate of convergence if $\pexp=1$ and $2\strong\step > \qexp$.
\end{itemize}
\end{restatable}

\beginrev
\begin{remark}
Getting an explicit estimate for the constant in the $\bigoh(\cdot)$ guarantee of \cref{thm:SOS} is quite involved but, up to logarithmic and subleading factors, Chung's lemma \cite{Chu54,Pol87} can be used to show that
\begin{enumerate*}
[\itshape a\upshape)]
\item
if $2\strong\step > \qexp$, it scales as $(\Const_{\bias}+\Const_{\sigma})/[(2\strong\step - \qexp)(1-\conf)]$ where $\Const_{\bias} = \sup_{\run} \curr[\step]\curr[\bbound]$ and $\Const_{\sigma} = \sup_{\run} \curr[\step]^2\curr[\noisevar]$;
\item
if $2\strong\step = \qexp$, it scales as %$(\step^{2} G^{2} + \step\Const_{\sigma}) / [(\strong\step - 1)(1-\conf)]$.
$(\Const_{\bias}+\Const_{\sigma})(1+ \max\{(2\strong\step)^2,4\strong\step\})/(1-\delta)$; 
and  
\item if $2\strong\step < \qexp$ as $(\Const_{\bias}+\Const_{\sigma})(1+ \max\{(2\strong\step)^2,4\strong\step\})/[(\qexp -2\strong\step)(1-\delta)]$.
\end{enumerate*}
\end{remark}
\endedit

Besides providing
%the conditions that need to be satisfied under \eqref{eq:SOS} in the
a general framework for achieving trajectory convergence, \cref{thm:SOS} gives the rates of convergence of the sequence of play to the Nash policy in question.
In particular, with this result in hand, one can confidently argue about the distance of the iterates of \eqref{eq:PG} from equilibrium in a series of different environments.
More to the point, this convergence guarantee allows the algorithm designer to adapt the parameters of the learning process according to the complexity and limitations of the environment, a feature which further highlights the significance of this result.

We should also note the delicate interplay between the method's step-size and the achieved convergence rate.
In the case of \cref{mod:full}, \cref{cor:SOS} suggests a step-size of the form $\curr[\step] = \Theta(1/\run)$, leading to a $\bigoh(1/\run)$ convergence rate.
As we show in the appendix, this rate can be improved:
in the deterministic case with perfect gradient information, \eqref{eq:PG} with a suitably chosen constant step-size achieves a \emph{geometric} convergence rate, \ie $\norm{\curr - \eq} = \bigoh(\exp(-\rho\run))$ for some $\rho>0$ (\cf \cref{prop:geometric} in \cref{app:SOS}).
By contrast, in the case of \cref{mod:stoch}, the $\bigoh(1/\run)$ rate we provide cannot be improved, even if the quadratic minorant \eqref{eq:strong} that characterizes \ac{SOS} policies holds \emph{globally} \textendash\ and this because the learning process is running against standard lower bounds from convex optimization \cite{Nes04,Bub15}.

Perhaps the most significant guarantee from a practical point of view is the $\bigoh(1/\sqrt{\run})$ convergence rate attained in \cref{mod:value} (\cf \cref{alg:reinforce,alg:greedy}).
This guarantee amounts to a $\bigoh(1/\run^{1/4})$ convergence rate in terms of the (non-squared) distance to equilibrium which, mutatis mutandis, represents a notable improvement over the $\bigoh(1/\run^{1/6})$ guarantee of \citet{LOPP22} (expressed in norm values).
Of course, the latter guarantee is global \textendash\ because the focus of \cite{LOPP22} is stochastic \emph{potential} games \textendash\ but it also concerns the ``best iterate'' of the process (not its ``last iterate''), so the two results are not immediately comparable.
However, a useful ``best-of-both-worlds'' heuristic that can be inferred by the combination of these works is that, given a budget of training episodes, \cref{alg:greedy} can be run with a constant step-size as per \cite{LOPP22} for a sufficient fraction of this budget, and then with a $\bigoh(1/\run)$ ``cooldown'' schedule for the rest.
In this way, after an aggressive ``exploration'' phase, the algorithm's $\bigoh(1/\run^{1/4})$ rate would kick in and supply faster stabilization to an \ac{SOS} policy.
%This would allow the algorithm's relatively fast capture by the basin of attraction of a \acl{SOS} policy, in which case the vanishing step-size schedule would achieve faster stabilization.



%----------------------------------------------------------------------
%% Strict
%----------------------------------------------------------------------
\subsection{Convergence to deterministic Nash policies}

Our last series of results concerns the rate of convergence to deterministic Nash policies in generic stochastic games.
As we discussed in \cref{sec:preliminaries},
%if $\eq$ is a deterministic Nash policy, it satisfies the inequality $\braket{\gvalue(\policy)}{\policy - \eq} \leq -\strong \norm{\policy - \eq}$ so, up to a change of coefficient, it also satisfies \eqref{eq:SOS}.
deterministic Nash policies also satisfy  \eqref{eq:SOS}, so the rate of convergence of \eqref{eq:PG} to such policies can be harvested directly from \cref{thm:SOS}.
However, as we show below, a simple projection tweak in \eqref{eq:SOS} can improve this rate dramatically.

The tweak in question is inspired by the geometry of $\policies$ around a deterministic policy:
by definition, such policies are corner points of $\policies$, so any consistent drift towards them will cause $\curr$ to hit the boundary of $\policies$ in finite time.
Of course, under \eqref{eq:PG}, the process may rebound from the boundary and return to the interior of $\policies$ if the policy gradient estimate is not particularly good at a given iteration of the algorithm.
However, if we replace the projection step of \eqref{eq:PG} with a ``lazy projection'' in the spirit of \citet{Zin03}, the aggregation of gradient steps will eventually push the process far inside the normal cone of $\policies$ at $\eq$, so rebounds of this type can no longer occur.

Formally, we will consider the following \acdef{LPG} scheme:
\begin{equation}
\label{eq:LPG}
\tag{\ac{LPG}}
\next[\dstate]
	= \curr[\dstate] + \curr[\step]\curr[\signal]
	\qquad
\next[\policy]
	= \proj_{\policies}(\next[\dstate])
\end{equation}
where $\curr[\dstate] = (\dstate_{\play,\run})_{\play\in\players} \in \dspace$ is an auxiliary variable that maintains an aggregate of gradient steps \emph{before} projecting them back to $\policies$.
%to obtain a new policy and continue playing.
We then have the following convergence result:

\begin{restatable}{theorem}{strict}
\label{thm:strict}
Let $\curr$ be the sequence of play under \eqref{eq:LPG} with step-size and policy gradient estimates such that $\pexp + \bexp >1$ and $\pexp - \sexp > 1/2$ as per \eqref{eq:schedule}.
If $\eq$ is a deterministic Nash policy, there exists an unbounded open set $\dbasin\subseteq\dspace$ of initializations such that, for any $\conf>0$, we have
\begin{equation}
\label{eq:lazy}
\probof{\text{$\curr$ converges to $\eq$} \given \init[\dstate] \in \dbasin}
	\geq 1-\conf,
\end{equation}
provided that $\step>0$ is small enough.
Moreover, conditioned on this event, $\curr$ converges to $\eq$ at a finite number of iterations, \ie there exists some $\run_{0}$ such that $\curr = \eq$ for all $\run\geq\run_{0}$.
\end{restatable}

\begin{corollary}
\label{cor:strict}
Suppose that \crefrange{mod:full}{mod:value} are run with parameters $\curr[\step] = \step/\run^{\pexp}$, $\pexp\in(1/2,1]$, and if applicable, $\curr[\mix] = \mix/\run^{\rexp}$ with $1-\pexp < \rexp < 2\pexp-1$.
Then the conclusions of \cref{thm:strict} hold.
\end{corollary}

\beginrev
\begin{remark}
Getting an explicit bound for $\run_{0}$ is quite complicated, but the last part of the proof of \cref{thm:strict} shows that $\run_{0}$ scales in terms of the parameters of the game and the algorithm as $n_0 = \bigoh\parens[\Big]{\parens[\big]{\frac{M\nStates\nPures}{\const\step}}^{1/(1-\pexp)}}$ where $\const>0$ measures the minimum payoff difference between equilibrium and non-equilibrium strategies at $\eq$, $M$ is a measure of the initial distance from $\eq$, and $\nStates$ and $\nPures$ is the number of states and pure strategies respectively.
\end{remark}
\endedit

\Cref{thm:strict} \textendash\ and, by extension, \cref{cor:strict} \textendash\ are fairly unique because they provide a guarantee for convergence to an \emph{exact} \acl{NE} in a \emph{finite} number of iterations.
To the best of our knowledge, the only comparable result in the literature is that of \cite{zhang2021gradient}, where the authors provide a finite-time convergence guarantee to strict equilibria with \emph{perfect} policy gradients (as per \cref{mod:full}).
The result of \citet{zhang2021gradient} echoes the convergence properties of deterministic first-order algorithms around sharp minima of convex functions \cite{Pol87}, but the fact that \cref{thm:strict} applies to models with \emph{stochastic} gradient feedback of \emph{unbounded} variance (\cref{mod:stoch,mod:value} respectively) is a major difference.
As far as we are aware, this is the first guarantee of its kind in the literature on learning in stochastic games.